\section{Screening vs.~Forensic Stages: An Architectural Test}
\label{sec:forensic}

The screen we have built so far runs on award records. The claim
of this section is that doing so is not a concession to data
scarcity---it is the right division of labour between the
screening stage and the forensic stage that follows.

Under incomplete operational observability of the bid layer
(\S\ref{sec:institutional}), enforcement architecture splits
naturally into two stages: an award-layer screening stage that
produces candidate flags, and a bid-layer forensic stage that
adjudicates suspected coordination on the bid distribution. We
test that division along three dimensions. The two stages
discriminate the same cartel-adjacency labels at comparable AUC
on different observability layers. The combination beats either
component alone. And the award-layer screen functions as a
Stage-1 gatekeeper, narrowing forensic-recoverable bid-microdata
interrogation to a manageable subpool while preserving most of
the joint-scoring recall.

Figure~\ref{fig:coarsening} sets up the comparison. The bid
layer carries the seven within-tender features the
Imhof--Wallimann pipeline aggregates into a discrimination score
(within-tender coefficient of variation, its standard deviation
across tenders, skewness, kurtosis, normalized spread, log
min/max, log-second-lowest distance); the operational layer
preserves only the award envelope---winner identity,
participant identity, item code, negotiated price---on which the
screening statistic operates. The screen the title refers to is
the award-layer rule alone; the bid-microdata pipeline appears
in this section as the benchmark we read the screen against and
as the Stage-2 input the screen triages, never as a screening
input itself.

\begin{figure}[htbp]
\centering
\includegraphics[width=0.92\textwidth]{./output/figures/fig_data_coarsening}
\caption[What survives data coarsening]{Bid-layer features and
award-layer fields. The bid-distribution features the
Imhof--Wallimann pipeline operates on (left, dashed arrows) are
dropped by the coarsening that defines most enforcement
environments at the operational layer; the four award-record
fields (right, solid arrows) survive. Boxes on the right report
random-forest AUCs against $\valCobidders$ adjudicated
CADE-cobidder labels (Table~\ref{tab:imhof_full}).}
\label{fig:coarsening}
\end{figure}

On the full always-loser pool ($N=\valImhofPoolN$ firms,
$\valCobidders$ adjudicated cobidder labels;
Table~\ref{tab:imhof_full}), the screening statistic
discriminates at AUC $\valImhofFLBin$ ($95\%$ CI
$\valImhofFLBinCI$) under the binary FL14 indicator and
$\valImhofFLcont$ ($\valImhofFLcontCI$) under the continuous
$\log(1+\text{tenders\_count})$ score. The full Imhof--Wallimann
pipeline reaches AUC $\valImhofFull$ ($\valImhofFullCI$) on the
same target. The two pipelines reach statistically
indistinguishable discrimination on different data envelopes;
the construct does not outperform the bid-distribution pipeline,
and we do not claim it does---the architectural point is that
comparable accuracy is achievable on a substantially thinner
data envelope. FL14 combined with the full Imhof pipeline
reaches AUC $\valImhofComboBin$ ($\valImhofComboBinCI$);
continuous loss intensity combined with the full pipeline
reaches $\valImhofComboCont$ ($\valImhofComboContCI$); the
combined classifiers gain $5$--$7$ AUC points over either single
construct. Restricting attention to the same-sample audit
(Table~\ref{tab:imhof_incremental}), where all detectors
evaluate the identical subset of firms---the subset for which
Imhof features can be computed because bid microdata are
available---the screening statistic adds $+\valImhofIncFL$ AUC
over the Imhof full pipeline alone (DeLong
$p = \valImhofIncFLp$) and the combined model adds
$\valImhofIncCombo$ AUC over Imhof full alone ($p < 0.001$). The
award-layer screen and the bid-layer forensic moments carry
non-redundant information about the same target. The subset
restriction is for the comparison, not for the screen: the screen
itself runs on the four award-record fields whether or not bid
microdata are recoverable for any given firm.

% comparison replaces the v13 strawman (CV-only) with the actual five-feature
% pipeline. Source: scripts/31_imhof_full_pipeline.R.
\begin{table}[htbp]
\hfuzz=300pt\relax
\centering
\caption{Full Imhof--Wallimann Pipeline vs Loser-Side Concentration: Discrimination at Different Data Costs}
\label{tab:imhof_full}
\resizebox{\textwidth}{!}{%
\begin{threeparttable}
\small
\begin{tabular}{lccc}
\toprule
Model & Features required & AUC & 95\% CI \\
\midrule
\multicolumn{4}{l}{\emph{Participation-only (award-record level)}} \\
FL flag (binary)            & participation count $>$ 14, win flag                   & $0.903$ & $[0.884, 0.923]$ \\
$\log(1{+}\text{tenders\_count})$ & participation count, win flag                    & $0.884$ & $[0.860, 0.908]$ \\
\midrule
\multicolumn{4}{l}{\emph{Bid-distribution (per-bid microdata)}} \\
Imhof CV only (v13 strawman) & per-bid offer values                                  & $0.585$ & $[0.553, 0.616]$ \\
Imhof full pipeline          & 5 within-tender features (CV, skew, kurt, spread, second-low) & $0.888$ & $[0.865, 0.911]$ \\
\midrule
\multicolumn{4}{l}{\emph{Combined}} \\
FL flag $+$ Imhof full       & participation count + 5 bid features                  & $0.955$ & $[0.943, 0.967]$ \\
$\log(1{+}\text{tenders\_count})$ $+$ Imhof full & participation count + 5 bid features & $0.962$ & $[0.954, 0.969]$ \\
\bottomrule
\end{tabular}
\begin{tablenotes}
\small
\item \textit{Notes:} 5-fold cross-validated AUC against $193$
CADE-cobidder labels within the always-loser pool ($N\!=\!16{,}779$).
Random forests with $500$ trees, otherwise default tuning. Bid-level
features are aggregated to the firm level by averaging the per-tender
within-tender statistic across each firm's tenders.
The two top blocks compare alternative single-construct screens
operating on different data envelopes; the bottom block reports the
combined classifiers. The combined model gains $5$--$7$ AUC points
over either single construct, indicating non-redundant signal in the
participation-only and bid-distribution information sources. The
participation-only screen achieves discrimination comparable to the
full Imhof--Wallimann pipeline at substantially lower data cost,
which is the operational claim of the paper. We do not claim the
participation-only screen \emph{outperforms} the bid-distribution
pipeline; the relevant comparison is at fixed data cost, where the
two pipelines reach similar accuracy and the combination dominates.
\end{tablenotes}
\end{threeparttable}
}
\end{table}


Beyond the AUC headline lies the operational test. The
award-layer screen functions as a Stage-1 gatekeeper for the
bid-layer forensic stage, narrowing the pool of firms whose bid
microdata the forensic stage needs to interrogate. The screen
itself does not need bid microdata; the gatekeeper claim is
about how it shapes the data envelope that Stage-2 must request.
Table~\ref{tab:architecture_gatekeeper} reports precision,
recall, and bid-microdata footprint for four rules on the same
evaluation pool of $\valArchPoolSize$ always-loser firms
($\valArchNPos$ adjudicated cobidder positives, base rate
$\valArchBaseRate$): the award-layer screen alone, the bid-layer
Imhof full pipeline alone, joint single-model scoring, and the
sequential gatekeeper that uses the award-layer screen to select
a Stage-1 pool of $K_1$ firms before applying Imhof full within
that pool. At top-$1{,}000$
flags, joint single-model scoring captures
$\valArchTPThouJoint$ of $\valArchNPos$ cobidders (recall
$\valArchRecallThouJoint$) while requiring bid microdata for the
entire $\valArchPoolSize$-firm pool; the sequential gatekeeper
with $K_1=\valArchSeqKone$ captures $\valArchTPThouSeqK$
cobidders (recall $\valArchRecallThouSeqK$) while requiring bid
microdata for $\valArchSeqKone$ firms only---a
$\valArchFootprintReduction$ reduction in bid-microdata
footprint at a $\valArchRecallCost$ recall cost relative to
joint scoring. The sequential rule is not the precision
champion (joint scoring is, at $\valArchPrecAtFiveHundJoint$
vs.\ $\valArchPrecAtFiveHundSeqK$ at top-$500$); its in-sample
operational case is the data-envelope wedge.

% Subprompt 3: architecture / gatekeeper (script 63)
\begin{table}[htbp]
\centering
\caption{Architecture / Gatekeeper: Operational Precision and Bid-Microdata Footprint by Rule}
\label{tab:architecture_gatekeeper}
\begin{threeparttable}
\footnotesize
\begin{tabular}{lrrrrrr}
\toprule
Rule & $k$ & TP & Precision & Recall & Lift & Bid-microdata firms needed \\
\midrule
Award-layer only (FL log\_tc) & $250$ & $40$ & $0.160$ & $0.207$ & $9.68\times$ & $0$ \\
 & $500$ & $66$ & $0.132$ & $0.342$ & $7.99\times$ & $0$ \\
 & $1{,}000$ & $97$ & $0.097$ & $0.503$ & $5.87\times$ & $0$ \\
\midrule
Bid-layer only (Imhof full) & $250$ & $50$ & $0.200$ & $0.259$ & $12.10\times$ & $11{,}676$ \\
 & $500$ & $73$ & $0.146$ & $0.378$ & $8.83\times$ & $11{,}676$ \\
 & $1{,}000$ & $97$ & $0.097$ & $0.503$ & $5.87\times$ & $11{,}676$ \\
\midrule
Joint scoring (FL + Imhof, single model) & $250$ & $74$ & $0.296$ & $0.383$ & $17.91\times$ & $11{,}676$ \\
 & $500$ & $112$ & $0.224$ & $0.580$ & $13.55\times$ & $11{,}676$ \\
 & $1{,}000$ & $142$ & $0.142$ & $0.736$ & $8.59\times$ & $11{,}676$ \\
\midrule
Sequential: FL $\rightarrow$ Imhof (Stage-1 keeps top 1000) & $250$ & $62$ & $0.248$ & $0.321$ & $15.00\times$ & $1{,}000$ \\
 & $500$ & $87$ & $0.174$ & $0.451$ & $10.53\times$ & $1{,}000$ \\
 & $1{,}000$ & $97$ & $0.097$ & $0.503$ & $5.87\times$ & $1{,}000$ \\
\midrule
Sequential: FL $\rightarrow$ Imhof (Stage-1 keeps top 2000) & $250$ & $65$ & $0.260$ & $0.337$ & $15.73\times$ & $2{,}000$ \\
 & $500$ & $96$ & $0.192$ & $0.497$ & $11.62\times$ & $2{,}000$ \\
 & $1{,}000$ & $131$ & $0.131$ & $0.679$ & $7.93\times$ & $2{,}000$ \\
\midrule
Sequential: FL $\rightarrow$ Imhof (Stage-1 keeps top 4000) & $250$ & $59$ & $0.236$ & $0.306$ & $14.28\times$ & $4{,}000$ \\
 & $500$ & $87$ & $0.174$ & $0.451$ & $10.53\times$ & $4{,}000$ \\
 & $1{,}000$ & $127$ & $0.127$ & $0.658$ & $7.68\times$ & $4{,}000$ \\
\addlinespace
\multicolumn{7}{l}{\textit{Reference: evaluation pool size = $11{,}676$ always-loser firms with $\geq 2$ priced bids per tender; $193$ cobidder positives; base rate $= 0.0165$.}} \\
\bottomrule
\end{tabular}
\begin{tablenotes}
\footnotesize
\item \textit{Notes:} Always-loser firms in the BEC sample, restricted to those for whom the within-tender bid-distribution moments (CV, SD, skewness, kurtosis, spread, min/max log, second-lowest distance --- the seven Imhof full-pipeline features) can be computed from \texttt{bid\_level\_with\_prices.parquet}. Five-fold cross-validated random forest scores; no out-of-sample temporal holdout in this table (see Online Appendix~C for the temporal-holdout audit). The single-model joint scoring trains a random forest on $\log(1+\text{tenders\_count})$ together with the seven Imhof features in one estimator; the sequential gatekeeper trains separately, using FL-stage rank to filter, then Imhof-stage rank inside the survivor pool. ``Bid-microdata firms needed'' is the number of always-loser firms for which per-bid \textit{Valor Unit\'ario Proposta} must be retrieved to compute the Imhof features under each rule. The award-layer-only rule needs zero bid-microdata extractions because $\log(1+\text{tenders\_count})$ uses participation counts only. Lift is precision-at-$k$ divided by the unconditional cobidder rate in the evaluation pool. The sequential pipeline does not raise top-$k$ precision above joint scoring; its operational claim is that comparable precision is reachable while requesting bid microdata for only the Stage-1 survivor pool, not the full pool.
\item \textit{Source:} \texttt{scripts/63\_architecture\_gatekeeper.R}.
\end{tablenotes}
\end{threeparttable}
\end{table}


In-sample evaluation invites the natural worry that the score
and the labels both draw on the full $2009$--$2019$ window: the
5-fold cross-validated precision-at-$k$ in
Table~\ref{tab:architecture_gatekeeper} confounds prediction
with classification. Table~\ref{tab:architecture_gatekeeper_th}
restricts every per-firm feature to participation observed in
$2009$--$2016$, leaving the cobidder ground truth as anchored on
post-sample CADE adjudications (which cannot be temporally
truncated without discarding nearly all positives); the
temporal-holdout pool is $\valArchTHPoolSize$ always-losers with
both train-window participation $>0$ and complete train-window
Imhof features ($\valArchTHNPos$ cobidder positives). Two
patterns matter for the architectural reading. Precision
compresses asymmetrically across the four rules: the award-layer
screen alone loses $\valArchTHFLDrop$ at top-$500$ (consistent
with the precision-at-$k$ audit on FL alone in Online
Appendix~C), joint scoring loses $\valArchTHJointDrop$, and the
sequential gatekeeper loses $\valArchTHSeqDrop$. In absolute
true-positive count, the sequential gatekeeper matches or
exceeds joint scoring under temporal holdout: at top-$500$ the
sequential rule recovers $\valArchTHTPFiveHundSeqK$ adjudicated
cobidders versus $\valArchTHTPFiveHundJoint$ for joint scoring,
and at top-$1{,}000$ it recovers $\valArchTHTPThouSeqK$ versus
$\valArchTHTPThouJoint$. Stage-2 Imhof moments derived from
train-window bid distributions are more temporally stable than
the FL participation rank applied across the full pool, so
restricting bid-microdata interrogation to a Stage-1 survivor
pool both compresses the data envelope and preserves out-of-time
TP capture better than running joint scoring on the full pool.
The $\valArchFootprintReduction$ data-envelope reduction does
not trade off against out-of-time recall.

% Subprompt 4 follow-on: temporal-holdout audit of the gatekeeper (script 64)
\begin{table}[htbp]
\centering
\caption{Architecture / Gatekeeper Under Temporal Holdout: Features Trained on $2009$--$2016$, Evaluated on the Same Cobidder Labels}
\label{tab:architecture_gatekeeper_th}
\begin{adjustbox}{max width=\textwidth,center}
\begin{threeparttable}
\footnotesize
\begin{tabular}{lrrrrrrr}
\toprule
 & & \multicolumn{2}{c}{Precision @ $k$} & \multicolumn{2}{c}{Recall @ $k$} & \multicolumn{2}{c}{TP @ $k$} \\
\cmidrule(lr){3-4} \cmidrule(lr){5-6} \cmidrule(lr){7-8}
Rule & $k$ & In-samp.\ & T.\ holdout & In-samp.\ & T.\ holdout & In-samp.\ & T.\ holdout \\
\midrule
Award-layer only & $250$ & $0.160$ & $0.076$ & $0.207$ & $0.134$ & $40$ & $19$ \\
 & $500$ & $0.132$ & $0.070$ & $0.342$ & $0.246$ & $66$ & $35$ \\
 & $1{,}000$ & $0.097$ & $0.066$ & $0.503$ & $0.465$ & $97$ & $66$ \\
Bid-layer only (Imhof full) & $250$ & $0.200$ & $0.156$ & $0.259$ & $0.275$ & $50$ & $39$ \\
 & $500$ & $0.146$ & $0.112$ & $0.378$ & $0.394$ & $73$ & $56$ \\
 & $1{,}000$ & $0.097$ & $0.082$ & $0.503$ & $0.577$ & $97$ & $82$ \\
Joint scoring (FL + Imhof) & $250$ & $0.296$ & $0.244$ & $0.383$ & $0.430$ & $74$ & $61$ \\
 & $500$ & $0.224$ & $0.170$ & $0.580$ & $0.599$ & $112$ & $85$ \\
 & $1{,}000$ & $0.142$ & $0.111$ & $0.736$ & $0.782$ & $142$ & $111$ \\
Sequential FL $\rightarrow$ Imhof, $K_1=2{,}000$ & $250$ & $0.260$ & $0.248$ & $0.337$ & $0.437$ & $65$ & $62$ \\
 & $500$ & $0.192$ & $0.174$ & $0.497$ & $0.613$ & $96$ & $87$ \\
 & $1{,}000$ & $0.131$ & $0.114$ & $0.679$ & $0.803$ & $131$ & $114$ \\
\bottomrule
\end{tabular}
\begin{tablenotes}
\footnotesize
\item \textit{Notes:} Temporal-holdout column restricts every per-firm feature to participation observed in $2009$--$2016$ before random-forest scoring. The cobidder ground truth (193 always-losers co-bidding with direct CADE defendants) is anchored on adjudications post-sample and cannot be temporally truncated without discarding nearly all positives, so the holdout applies to features rather than labels. Evaluation pool: in-sample is the $11{,}676$-firm pool of Table~\ref{tab:architecture_gatekeeper}; temporal-holdout pool is the $8{,}257$ always-losers with both train-window participation $>0$ and complete train-window Imhof features. The pattern parallels the precision-at-$k$ audit on FL alone reported in Online Appendix C: precision compresses by roughly half, recall by less, lift by roughly half. The relative ordering of the four rules is preserved.
\item \textit{Source:} \texttt{scripts/64\_gatekeeper\_temporal\_holdout.R}.
\end{tablenotes}
\end{threeparttable}
\end{adjustbox}
\end{table}


Taken together, the three findings identify the sequencing
structure the paper proposes. Under incomplete operational
observability, the deployable architecture is an award-layer
screening stage running on operational data, followed by a
bid-layer forensic stage running the bid-distribution moment
battery on the candidates the screening stage flags, with bid
microdata recovered through administrative request only for the
screening-flagged subpool. Jurisdictions whose operational layer
carries award records but not per-bidder bid amounts can deploy
the screening stage immediately. Jurisdictions whose
forensic-recoverable layer is also analytically queryable gain
the forensic stage on top, with the screen already constraining
the pool of cases that warrant the forensic step. Reforms that
mandate operational bid-microdata archival
(\S\ref{sec:institutional}) are best read as expanding the
forensic stage rather than substituting for the screening one---a
reading that reverses the standard policy intuition that
operational bid-archival is a precondition for proactive cartel
screening.
